% ------------------------------------------------------------------
%  Linkage criteria: single, complete, average, Ward.
%
%  Not written from method_template.tex: these are criteria shared by
%  every hierarchical method, not methods themselves. The concrete
%  hierarchical methods (Ward in Task 40 first) will each get a section
%  from method_template that names the criterion it fits under.
%
%  Code counterpart: xxcluster/cluster/hierarchical/linkage.py.
%
%  NOTE: citation keys marked ref_TODO below need replacement from the
%  shared citation-key mapping sheet before merge. Lance-Williams
%  (1967) and Murtagh-Legendre (2014) are not yet in the bibliography.
%  ref_11 (Han/Kamber/Pei) is already present and used in-place.
% ------------------------------------------------------------------
\subsubsection{Linkage criteria}
\label{sec:tech:linkage}

\paragraph{Overview}
\label{sec:tech:linkage:overview}
A linkage criterion lifts a dissimilarity between observations to a
dissimilarity between clusters, and is the component every hierarchical
method of \textbf{Sect.~\ref{sec:background:families}} shares. The
criterion is kept separate from the method for two reasons: the same
criterion is reused across methods, and the choice of criterion changes
the shape of cluster a method can recover, which is a result to report
rather than an implementation detail~\cite{ref_11}. Four criteria are
implemented here (single, complete, average and Ward); the first three
apply to any dissimilarity and the fourth requires Euclidean distance.

\paragraph{Definition}
\label{sec:tech:linkage:definition}
Given a dissimilarity $d(\cdot, \cdot)$ on observations and two disjoint
clusters $\mathcal C_A, \mathcal C_B$, each criterion returns a
non-negative $d(\mathcal C_A, \mathcal C_B)$:
\begin{align}
    \text{Single:}   \quad & d(\mathcal C_A, \mathcal C_B)
        = \min_{\mathbf a \in \mathcal C_A, \mathbf b \in \mathcal C_B}
          d(\mathbf a, \mathbf b), \\
    \text{Complete:} \quad & d(\mathcal C_A, \mathcal C_B)
        = \max_{\mathbf a \in \mathcal C_A, \mathbf b \in \mathcal C_B}
          d(\mathbf a, \mathbf b), \\
    \text{Average (UPGMA):} \quad & d(\mathcal C_A, \mathcal C_B)
        = \tfrac{1}{|\mathcal C_A| \, |\mathcal C_B|}
          \textstyle \sum_{\mathbf a \in \mathcal C_A}
                     \sum_{\mathbf b \in \mathcal C_B}
          d(\mathbf a, \mathbf b), \\
    \text{Ward:} \quad & d(\mathcal C_A, \mathcal C_B)
        = \tfrac{|\mathcal C_A| \, |\mathcal C_B|}
                {|\mathcal C_A| + |\mathcal C_B|}
          \left\| \boldsymbol \mu_A - \boldsymbol \mu_B \right\|_2^2,
\end{align}
where $\boldsymbol \mu_A, \boldsymbol \mu_B$ are the cluster centroids
and Ward's expression is the increase in the within-cluster
sum-of-squared-error upon merging \cite{ref_TODO_ward1963}.

\paragraph{Lance--Williams recurrence}
\label{sec:tech:linkage:recurrence}
Naively computing $d(\mathcal C_A, \mathcal C_B)$ from the definition
on every merge is $\mathcal O(m^2)$ at each of the $m - 1$ merges, and
therefore unusable at any real data volume. Lance and Williams observed
that all four criteria above (and several others) express the
dissimilarity between a newly merged cluster $\mathcal C_A \cup
\mathcal C_B$ and every other cluster $\mathcal C_I$ as a linear
combination of the three prior dissimilarities
\cite{ref_TODO_lance_williams}:
\begin{align}
    d(\mathcal C_A \cup \mathcal C_B, \mathcal C_I)
    = \alpha_A \, d(\mathcal C_A, \mathcal C_I)
    + \alpha_B \, d(\mathcal C_B, \mathcal C_I)
    + \beta   \, d(\mathcal C_A, \mathcal C_B)
    + \gamma  \, \left| d(\mathcal C_A, \mathcal C_I) - d(\mathcal C_B, \mathcal C_I) \right|.
\end{align}
The coefficients $(\alpha_A, \alpha_B, \beta, \gamma)$ are per criterion
and, for Ward, depend on the cluster sizes:
\begin{center}
\begin{tabular}{lcccc}
\toprule
Criterion & $\alpha_A$ & $\alpha_B$ & $\beta$ & $\gamma$ \\
\midrule
Single   & $\tfrac{1}{2}$ & $\tfrac{1}{2}$ & $0$ & $-\tfrac{1}{2}$ \\
Complete & $\tfrac{1}{2}$ & $\tfrac{1}{2}$ & $0$ & $+\tfrac{1}{2}$ \\
Average (UPGMA) & $\tfrac{|\mathcal C_A|}{|\mathcal C_A|+|\mathcal C_B|}$
                & $\tfrac{|\mathcal C_B|}{|\mathcal C_A|+|\mathcal C_B|}$
                & $0$ & $0$ \\
Ward     & $\tfrac{|\mathcal C_A|+|\mathcal C_I|}{|\mathcal C_A|+|\mathcal C_B|+|\mathcal C_I|}$
         & $\tfrac{|\mathcal C_B|+|\mathcal C_I|}{|\mathcal C_A|+|\mathcal C_B|+|\mathcal C_I|}$
         & $-\tfrac{|\mathcal C_I|}{|\mathcal C_A|+|\mathcal C_B|+|\mathcal C_I|}$
         & $0$ \\
\bottomrule
\end{tabular}
\end{center}
Overriding \texttt{update} with this recurrence, rather than falling
back to the $\mathcal O(m^2)$ \texttt{between}, is what makes the merge
loop affordable and is a requirement rather than an optimisation for
this project's data volume.

\paragraph{Regime shapes recovered}
\label{sec:tech:linkage:shapes}
The four criteria are not interchangeable: each imposes a shape bias
that determines the regime a hierarchical method can recover under it,
and getting this wrong is how a real operating structure becomes an
artefact of the criterion.
\begin{description}
    \item[Single] chains observations. A single bridging point between
    two otherwise distinct groups is enough to merge them, since only
    the nearest pair of members is consulted. This is useful for
    genuinely elongated regimes (a slow drift along one operating
    variable), and pathological on dense plant data where a stray
    reading between two envelopes joins them into one meaningless
    regime~\cite{ref_11}.

    \item[Complete] enforces compactness. A merge is charged by its
    worst pairwise disagreement, so clusters remain tight and
    comparably sized, and the criterion refuses to absorb an outlier
    into either group. The corresponding failure is splitting a real
    elongated regime into pieces, one per part of it that happens to
    lie within its diameter~\cite{ref_11}.

    \item[Average (UPGMA)] is the compromise. Less prone to chaining
    than single because it consults every pair; less brittle to
    boundary observations than complete because no single pair
    dominates. It has no strong shape bias of its own and is the
    default reach when the regime shape is not known in advance.

    \item[Ward] minimises within-cluster variance. It rewards spherical,
    equally-sized clusters -- the same bias as $k$-means, since the
    criterion is derived from the same objective -- and therefore
    combines with a validity index that shares the bias
    (Sect.~\ref{sec:measure:silhouette}) to give a compact
    representation of the operating record. The bias makes it the
    wrong choice where regimes are elongated or unequal in size, and
    it is undefined under any non-Euclidean dissimilarity for the
    reason given below.
\end{description}

\paragraph{Ward's Euclidean requirement}
\label{sec:tech:linkage:ward}
The Ward increase written in the definition uses squared Euclidean
distance between centroids, and the Lance--Williams recurrence above
tracks that quantity as the criterion. Under any other dissimilarity
the recurrence is not the criterion it names and a method built on it
would silently return a partition it did not compute. The
implementation therefore declares \texttt{requires\_euclidean = True},
and \texttt{BaseHierarchicalClusterer.\_fit} refuses a
non-Euclidean measure at fit time rather than at the boundary. Two
distinct definitions in the literature both call themselves Ward: the
implementation here is the ``Ward2'' variant matching
\texttt{scipy.cluster.hierarchy.linkage(method="ward")} and
scikit-learn's Agglomerative Ward
\cite{ref_TODO_murtagh_legendre}.

\paragraph{Monotonicity}
\label{sec:tech:linkage:monotonic}
A monotonic criterion produces a dendrogram with non-decreasing merge
heights, so a height threshold is a valid cut. All four criteria
implemented here are monotonic; each declares this as a class attribute
that \texttt{BaseHierarchicalClusterer.cut} reads as the cheap guard,
alongside its own check on the fitted heights. The two checks agree by
construction, and they must: the declared attribute is the guard that
refuses an inversion-producing sweep before it runs.

\paragraph{Verification}
\label{sec:tech:linkage:verification}
Each criterion is verified in \texttt{tests/test\_linkage.py} to
reproduce the merge heights of
\texttt{scipy.cluster.hierarchy.linkage(method=\ldots)} on the iris
benchmark, to floating-point tolerance. A separate fixture --- five
points arranged as two tight groups joined by a stepping stone ---
evidences the chains-versus-compacts distinction stated above, with
single linkage joining the two groups at the intermediate height and
complete linkage keeping them apart until the last merge.

\paragraph{Summary}
\label{sec:tech:linkage:summary}
Four registered criteria fix the shape of cluster the hierarchical
methods of Sect.~\ref{sec:tech:hierarchical} can recover. Each is
implemented with the Lance--Williams recurrence to keep the merge loop
affordable, each declares its Euclidean requirement and its
monotonicity, and each reproduces its SciPy counterpart in the test
suite. The first hierarchical method built on them, Ward's method, is
introduced in Sect.~\ref{sec:tech:ward} (Task 40).
